\documentclass[10pt,a4paper]{amsart}
\usepackage{amsmath,amssymb,amsthm,graphicx}
% For commutative diagrams we use Xy-pic
\usepackage[all]{xy}

% We use 2cell for 2-commutative diagrams.
\xyoption{2cell}
% Theorem environments.
%
\theoremstyle{plain}
\newtheorem{theorem}[subsection]{Theorem}
\newtheorem{proposition}[subsection]{Proposition}
\newtheorem{lemma}[subsection]{Lemma}

\theoremstyle{definition}
\newtheorem{definition}[subsection]{Definition}
\newtheorem{example}[subsection]{Example}
\newtheorem{exercise}[subsection]{Exercise}
\newtheorem{situation}[subsection]{Situation}

\theoremstyle{remark}
\newtheorem{remark}[subsection]{Remark}
\newtheorem{remarks}[subsection]{Remarks}

\numberwithin{equation}{subsection}

% Macros
%
\def\lim{\mathop{\mathrm{lim}}\nolimits}
\def\colim{\mathop{\mathrm{colim}}\nolimits}
\def\Spec{\mathop{\mathrm{Spec}}}
\def\Hom{\mathop{\mathrm{Hom}}\nolimits}
\def\Ext{\mathop{\mathrm{Ext}}\nolimits}
\def\SheafHom{\mathop{\mathcal{H}\!\mathit{om}}\nolimits}
\def\SheafExt{\mathop{\mathcal{E}\!\mathit{xt}}\nolimits}
\def\Sch{\mathit{Sch}}
\def\Mor{\mathop{\mathrm{Mor}}\nolimits}
\def\Ob{\mathop{\mathrm{Ob}}\nolimits}
\def\Sh{\mathop{\mathit{Sh}}\nolimits}
\def\NL{\mathop{N\!L}\nolimits}
\def\CH{\mathop{\mathrm{CH}}\nolimits}
\def\proetale{{pro\text{-}\acute{e}tale}}
\def\etale{{\acute{e}tale}}
\def\QCoh{\mathit{QCoh}}
\def\Ker{\mathop{\mathrm{Ker}}}
\def\Im{\mathop{\mathrm{Im}}}
\def\Coker{\mathop{\mathrm{Coker}}}
\def\Coim{\mathop{\mathrm{Coim}}}

% Boxtimes
%
\DeclareMathSymbol{\boxtimes}{\mathbin}{AMSa}{"02}

%
% Macros for moduli stacks/spaces
%
\def\QCohstack{\mathcal{QC}\!\mathit{oh}}
\def\Cohstack{\mathcal{C}\!\mathit{oh}}
\def\Spacesstack{\mathcal{S}\!\mathit{paces}}
\def\Quotfunctor{\mathrm{Quot}}
\def\Hilbfunctor{\mathrm{Hilb}}
\def\Curvesstack{\mathcal{C}\!\mathit{urves}}
\def\Polarizedstack{\mathcal{P}\!\mathit{olarized}}
\def\Complexesstack{\mathcal{C}\!\mathit{omplexes}}
% \Pic is the operator that assigns to X its picard group, usage \Pic(X)
% \Picardstack_{X/B} denotes the Picard stack of X over B
% \Picardfunctor_{X/B} denotes the Picard functor of X over B
\def\Pic{\mathop{\mathrm{Pic}}\nolimits}
\def\Picardstack{\mathcal{P}\!\mathit{ic}}
\def\Picardfunctor{\mathrm{Pic}}
\def\Deformationcategory{\mathcal{D}\!\mathit{ef}}
\usepackage{cleveref}
\crefname{originallemma}{original lemma}{original lemmas}
\Crefname{originallemma}{Original lemma}{Original lemmas}
\makeatletter
\newlabel{more-algebra-lemma-henselization-normal}{{06DI}{source}{}{sourcepointer-1}{}}
\newlabel{more-algebra-lemma-henselization-normal@cref}{{[originallemma][0][]06DI}{[1][0][]source}{}{}{}}
\makeatother
\crefname{originalreference}{original reference}{original references}
\Crefname{originalreference}{Original reference}{Original references}
\makeatletter
\newlabel{theorem-higher-direct-images}{{03Q9}{source}{}{sourcepointer-2}{}}
\newlabel{theorem-higher-direct-images@cref}{{[originalreference][0][]03Q9}{[1][0][]source}{}{}{}}
\makeatother
\title{On an integral normal scheme the pushforward of a constant sheaf along the generic point or a nonempty open is again constant}
\author{The Stacks Project Authors}
\date{stacks-project:etale-cohomology.tex:local-snapshot-20261011; GNU FDL 1.2 or later}
\begin{document}
\maketitle
\noindent\textit{Editorial note.} All delivered mathematical text is the author's own, copied from the pinned Stacks Project chapter file etale-cohomology.tex; no statement or proof step was rewritten, reordered or removed.\par\medskip
\noindent\textit{Editorial note.} The counted proof is the author's complete argument: the proof is written out in full for both parts: the normality of the strict henselization and the resulting irreducibility for the first, and the injectivity from the density of the generic fibre for the second.\par\medskip
\noindent\textit{Editorial note.} The label lemma-push-constant-sheaf-from-open names the counted statement, and the proof environment immediately following it in the source is delivered verbatim as its proof.\par\medskip
\noindent\textit{Editorial note.} The counted statement is carried by the author's own lemma environment in the chapter source.\par\medskip
\noindent\textit{Editorial note.} The statement and proof are copied from the pinned chapter file; the two cited inputs are kept as external pointers with their canonical Stacks tags.\par\medskip
\noindent\textit{Allowed standard prerequisites (editorial).}\begin{itemize}
\item Strict henselisations of a normal local ring are normal domains. More on Algebra, the lemma stating that a strict henselization of a normal local ring is a normal domain. Supplies the irreducibility of the base change Spec(K) x\_X Spec(O\textasciicircum{}sh\_\{X,x\}).
\item Stalks of a pushforward as cohomology of the fibre. Same chapter, the theorem describing the stalk of a pushforward as the cohomology of the corresponding fibre product. Supplies the identification of the stalk of g\_*underline E with E.
\item Density of the generic fibre over an etale neighbourhood. Standard property of an integral scheme, used in the second part: for V etale over X the scheme Spec(K) x\_X V is dense in U x\_X V for a nonempty open U. Supplies the injectivity of the map j\_*underline E -> g\_*underline E.
\end{itemize}\medskip
\ifdefined\hypertarget\hypertarget{sourcepointer-1}{}\fi
\noindent\textit{Source reference.} Original lemma 06DI: sources/stacks-chow/tags.txt, source lines 8182--8182. Prerequisite pointer invoked at proof line(s) 12715; it names a standard fact of a sibling chapter that the counted argument uses, and every step of the counted proof built on it is retained verbatim in the delivered body.\par\medskip
\ifdefined\hypertarget\hypertarget{sourcepointer-2}{}\fi
\noindent\textit{Source reference.} Original reference 03Q9: sources/stacks-etale-cohomology/pinned-original-source.tex, source lines 8015--8015. Prerequisite pointer invoked at proof line(s) 12720; it names a standard fact of the same chapter that the counted argument uses, and every step of the counted proof built on it is retained verbatim in the delivered body.\par\medskip
\begin{lemma}
\label{lemma-push-constant-sheaf-from-open}
Let $X$ be an integral normal scheme with function field $K$.
Let $E$ be a set.
\begin{enumerate}
\item Let $g : \Spec(K) \to X$ be the inclusion of the generic point.
Then $g_*\underline{E} = \underline{E}$.
\item Let $j : U \to X$ be the inclusion of a nonempty open. Then
$j_*\underline{E} = \underline{E}$.
\end{enumerate}
\end{lemma}

\begin{proof}
Proof of (1). Let $x \in X$ be a point. Let
$\mathcal{O}^{sh}_{X, \overline{x}}$
be a strict henselization of $\mathcal{O}_{X, x}$. 
By More on Algebra, Lemma \ref{more-algebra-lemma-henselization-normal}
we see that $\mathcal{O}^{sh}_{X, \overline{x}}$ is a normal domain.
Hence $\Spec(K) \times_X \Spec(\mathcal{O}^{sh}_{X, \overline{x}})$
is irreducible. It follows that
the stalk $(g_*\underline{E})_{\overline{x}}$ is equal to $E$,
see Theorem \ref{theorem-higher-direct-images}.

\medskip\noindent
Proof of (2). Since $g$ factors through $j$ there is a map
$j_*\underline{E} \to g_*\underline{E}$. This map is injective because
for every scheme $V$ \'etale over $X$ the set $\Spec(K) \times_X V$
is dense in $U \times_X V$. On the other hand, we have a map
$\underline{E} \to j_*\underline{E}$ and we conclude.
\end{proof}

\end{document}
